\section{Protocols and services}

\subsection{DAPHNE}
\label{sec:services-daphne}

\begin{center}\small
\begin{tabular}{L{0.25\textwidth}L{0.68\textwidth}}
\toprule
\textbf{Location} & \textbf{Service chain} \\
\midrule
DAPHNE & AXI / OS $\rightarrow$ \code{daphne-server} $\rightarrow$ ZMQ. \\
SC bridge server & ZMQ client $\rightarrow$ OPC UA bridge. \\
SC server & OPC UA client $\rightarrow$ Ignition. \\
DAQ server & DAQ CCM client/opmon; separate Hermes configuration client. \\
\bottomrule
\end{tabular}
\end{center}

The OPC UA bridge runs on the SC bridge server and includes DAQ configuration
and control arbitration. Its ZMQ client receives \code{daphne-server}
publications; the bridge exposes them to the SC server's OPC UA client and
to DAQ opmon. The SC server hosts Ignition. Hermes remains a separate DAPHNE
service receiving configuration over UDP/IPBus. The installation record shall
identify each service host and endpoint.

\begin{figure}[htbp]
\centering
\begin{tikzpicture}
\node[box,text width=37mm] (ignition) at (0,0)
 {SC server\\OPC UA client\\$\rightarrow$ Ignition};
\node[box,text width=32mm] (daq) at (5.0,0) {DAQ CCM client\\opmon subscriber};
\node[box,text width=35mm] (hermesclient) at (10.7,0)
 {DAQ Hermes\\configuration client};
\node[box,text width=75mm,minimum height=19mm] (ua) at (2.5,-2.3)
 {SC bridge server\\ZMQ client $\rightarrow$ OPC UA bridge\\Integrated DAQ configuration};
\node[box,text width=49mm,minimum height=18mm] (daphne) at (2.5,-5.8)
 {DAPHNE\\AXI / OS $\rightarrow$ \code{daphne-server}\\$\rightarrow$ ZMQ};
\node[box,text width=37mm,minimum height=18mm] (hermes) at (10.7,-5.8)
 {Board-side Hermes\\IPBus / UDP\\Configuration endpoint};
\node[draw=dunegreen,dashed,rounded corners,fit=(daphne)(hermes),inner sep=8pt,
 label={[font=\footnotesize,text=dunegreen]below:One DAPHNE board; one shared 1\,Gb Ethernet CCM connection}] {};
\draw[both] (ignition) -- node[left,font=\footnotesize]{OPC UA} (ignition |- ua.north);
\draw[both] (daq) -- node[right,font=\footnotesize]{OPC UA} (daq |- ua.north);
\draw[both] ([xshift=-11mm]ua.south) -- node[left,align=right,font=\footnotesize,fill=white,inner sep=2pt]
 {Configuration\\request/response} ([xshift=-11mm]daphne.north);
\draw[flow,draw=dunegreen] ([xshift=11mm]daphne.north) -- node[right,align=left,font=\footnotesize,fill=white,inner sep=2pt]
 {Published\\monitoring} ([xshift=11mm]ua.south);
\draw[both] (hermesclient.south) -- node[right,align=left,font=\footnotesize]
 {IPBus / UDP\\configuration} (hermes.north);
\end{tikzpicture}
\caption{DAPHNE service locations and protocol paths.}
\label{fig:protocol}
\end{figure}

\subsubsection{Configuration exchanges}

DAQ CCM requests run-variable and timing-endpoint configuration through the
OPC UA bridge. The ZMQ client on the SC bridge server exchanges Protobuf
messages over ZeroMQ and TCP using a DEALER client and the board's ROUTER
endpoint [R3, R4].
The reference port is TCP 40001 [R5]. Correlated responses shall distinguish
acceptance, completion and failure.

Hermes retains its independent IPBus/UDP configuration endpoint, reference
UDP 50001 [R6], outside the OPC UA bridge. Both services share the board
connection and gateware lifecycle.

\subsubsection{Published monitoring}

Monitoring shall be published periodically or on change, without a client
request for each update. Only configuration and control operations, including
reset/recovery commands, use application request/response. Local hardware
sampling is performed by the producing service.

The proposed DAPHNE publication transport is Protobuf over ZeroMQ PUB/SUB over TCP,
using endpoints separate from configuration. The ZMQ client on the SC bridge
server subscribes to \code{daphne-server} publications. The OPC UA bridge
distributes them through OPC UA subscriptions to the SC server and DAQ opmon.

\begin{center}\small
\begin{tabular}{L{0.25\textwidth}L{0.68\textwidth}}
\toprule
\textbf{Consumer} & \textbf{Monitoring responsibility} \\
\midrule
DAQ opmon & Firmware counters, timing state and applied data-readout configuration, published by \code{daphne-server} through the OPC UA bridge. \\
SC / Ignition & Voltages, temperatures, fans, service status and software/firmware versions. \\
SC and DAQ & Timing-endpoint condition, reset/recovery progress and board readiness. \\
\bottomrule
\end{tabular}
\end{center}

Opmon shall not acquire voltage monitoring. SC receives equipment monitoring
through the OPC UA bridge without an opmon dependency. Publications shall
carry source identity, time, sequence and validity. Subscribers shall detect
stale or missing updates and obtain a current snapshot after reconnecting.

\subsubsection{Timing-endpoint control and recovery over GbE}
\label{sec:timing-recovery}

Timing-endpoint configuration follows the DAQ configuration path:
\begin{center}
DAQ $\longleftrightarrow$ OPC UA bridge $\longleftrightarrow$
\code{daphne-server} $\longleftrightarrow$ local timing endpoint.
\end{center}
The same interface shall expose reset and recovery to authorized SC clients.
Timing monitoring shall be published over the board's GbE connection and
made available to both SC and DAQ.

A reset request shall identify its target and receive an execution response.
The board shall publish reset assertion, release and the resulting endpoint
state so that both clients can acknowledge the same recovery operation.
Recovery access shall remain available while timing is unavailable. Concurrent
DAQ and SC commands shall be serialized by SC's control infrastructure.

SC shall determine and publish \emph{board ready for data taking} from the
agreed readiness conditions, including completed reset, restored clocks,
endpoint readiness and valid synchronization [R4]. Reset release alone is
insufficient. Missing or stale required publications shall withhold readiness.
DAQ consumes the published readiness before resuming data taking.

\subsection{Light Calibration Module}
\label{sec:services-calibration}

\begin{center}\small
\begin{tabular}{L{0.25\textwidth}L{0.68\textwidth}}
\toprule
\textbf{Location} & \textbf{Calibration services} \\
\midrule
Light Calibration Module & PDS-provided LCM controller firmware and native calibration service. \\
SC physical server & Calibration OPC UA adapter and command arbitration. \\
SC infrastructure & Ignition client for equipment control and monitoring. \\
DAQ infrastructure & Client of the calibration interface for settings, sequences and status. \\
\bottomrule
\end{tabular}
\end{center}
SC provides equipment power/service control and monitoring. The installation
record shall identify each service host and endpoint.

Calibration requests and responses pass through the calibration OPC UA adapter
to the LCM native service. That service shall publish equipment health,
calibration status and readiness through OPC UA subscriptions to SC and DAQ.
The control interface shall expose protective inhibit actions.

PDS shall specify the native control protocol, transport, ports, command acknowledgements,
monitoring schema and compatible versions. These details remain open.

\subsubsection{Timing-endpoint control and recovery over GbE}

LCMs use the same timing protocol as DAPHNE. The requirements in
Section~\ref{sec:timing-recovery} apply to each LCM over its 1\,GbE CCM link,
through the LCM service and OPC UA adapter: DAQ configuration, authorized SC
reset/recovery, published acknowledgement and endpoint state, and command
arbitration. Recovery shall remain available during a timing outage.
SC shall publish LCM readiness from fresh timing and recovery evidence;
DAQ consumes it before calibration or data taking.

\subsection{PDS PoF laser boxes}
\label{sec:services-pof}

\begin{center}\small
\begin{tabular}{L{0.25\textwidth}L{0.68\textwidth}}
\toprule
\textbf{Location} & \textbf{PoF services} \\
\midrule
Compact PLC; placement to be specified & PoF control and monitoring logic. \\
Server location undecided & OPC UA server software for the compact PLC; implementation and driver to be selected. \\
SC infrastructure & Ignition OPC UA client for PoF control and monitoring. \\
DPS infrastructure & Detector protection; PLC/protection connection to be defined. \\
\bottomrule
\end{tabular}
\end{center}
SC and PDS shall design the compact PLC system and identify and test its
OPC UA server software. PDS provides PLC hardware and cables; SC integrates
control and monitoring into Ignition.

The PLC-to-OPC-UA protocol, driver, transport, ports and server location are
undecided. The selected implementation shall provide configuration/control
request/response and publish monitoring to Ignition through OPC UA
subscriptions. PoF has no DAQ interface.

The interface shall expose DPS protection status and give protective inhibits
priority over operating commands. PDS and DPS shall define the protection
connection and acknowledgement behavior. The interface release shall specify
the monitoring schema, compatible versions and service endpoints.
